\chapter{2007年湖北卷（理）}

\section{单选题}

\begin{problem}
如果 \( \left(3x^{2}-\frac{2}{x^{3}}\right)^{n} \) 的展开式中含有非零常数项，则正整数 \(n\) 的最小值为（\(\quad\)）

\choices
    {\(3\)}
    {\(5\)}
    {\(6\)}
    {\(10\)}
\end{problem}

\begin{answer}
B
\end{answer}

\begin{solution}
展开式的通项为 \(\mathrm C_n^k(3x^2)^{n-k}\left(-\dfrac{2}{x^3}\right)^k\)，其中 \(k=0\)，\(1\)，\(\ldots\)，\(n\)，其幂次为 \(2(n-k)-3k=2n-5k\). 含有非零常数项时，必须有 \(2n-5k=0\)，即 \(2n=5k\). 因为 \(n\) 为正整数，且 \(k\) 也为整数，所以 \(n\) 必须是 \(5\) 的倍数；取 \(n=5\) 时 \(k=2\) 满足 \(0\leqslant k\leqslant n\). 因而最小值为 \(5\)，故选 B.
\end{solution}

\begin{problem}
将 \( y = 2\cos \left(\frac{x}{3} + \frac{\pi}{6}\right) \) 的图象按向量 \( \bs{a} = \left(-\frac{\pi}{4}, -2\right) \) 平移，则平移后所得图象的解析式为（\(\quad\)）

\choices
    {\(y = 2\cos \left(\frac{x}{3} + \frac{\pi}{4}\right) - 2\)}
    {\(y = 2\cos \left(\frac{x}{3} - \frac{\pi}{4}\right) + 2\)}
    {\(y = 2\cos \left(\frac{x}{3} - \frac{\pi}{12}\right) - 2\)}
    {\(y = 2\cos \left(\frac{x}{3} + \frac{\pi}{12}\right) + 2\)}
\end{problem}

\begin{answer}
A
\end{answer}

\begin{solution}
设原函数为 \(f(x)=2\cos\left(\dfrac{x}{3}+\dfrac{\pi}{6}\right)\). 平移向量为 \(\left(-\dfrac{\pi}{4},-2\right)\) 时，原图象上的点 \((u,f(u))\) 变为 \(\left(u-\dfrac{\pi}{4},f(u)-2\right)\). 令新横坐标为 \(x=u-\dfrac{\pi}{4}\)，则 \(u=x+\dfrac{\pi}{4}\)，所以平移后的函数为
\[
y=f\left(x+\frac{\pi}{4}\right)-2=2\cos\left(\frac{x}{3}+\frac{\pi}{12}+\frac{\pi}{6}\right)-2=2\cos\left(\frac{x}{3}+\frac{\pi}{4}\right)-2
\]
故选 A.
\end{solution}

\begin{problem}
设 \(P\) 和 \(Q\) 是两个集合，定义集合 \(P - Q = \{x \mid x \in P, \text{且 } x \notin Q\}\)，如果 \(P = \{x \mid \log_2 x < 1\}\)，\(Q = \{x \mid |x - 2| < 1\}\)，那么 \(P - Q\) 等于（\(\quad\)）

\choices
    {\(\{x\mid 0 < x < 1\}\)}
    {\(\{x\mid 0 < x\leqslant 1\}\)}
    {\(\{x\mid 1\leqslant x < 2\}\)}
    {\(\{x\mid 2\leqslant x < 3\}\)}
\end{problem}

\begin{answer}
B
\end{answer}

\begin{solution}
由 \(\log_2x<1\) 且对数定义域为 \(x>0\)，得 \(P=(0,2)\). 由 \(|x-2|<1\) 得 \(1<x<3\)，即 \(Q=(1,3)\). 因而从 \(P\) 中去掉同时属于 \(Q\) 的元素，得到 \(P-Q=(0,1]\)，故选 B.
\end{solution}

\begin{problem}
平面 \(\alpha\) 外有两条直线 \(m\) 和 \(n\)，如果 \(m\) 和 \(n\) 在平面 \(\alpha\) 内的射影分别是 \(m'\) 和 \(n'\)，给出下列四个命题：

\begin{circlelist}
\item \(m' \perp n' \Rightarrow m \perp n\)；
\item \(m \perp n \Rightarrow m' \perp n'\)；
\item \(m'\) 与 \(n'\) 相交 \(\Rightarrow m\) 与 \(n\) 相交或重合；
\item \(m'\) 与 \(n'\) 平行 \(\Rightarrow m\) 与 \(n\) 平行或重合.
\end{circlelist}

其中不正确的命题个数是（\(\quad\)）

\choices
    {\(1\)}
    {\(2\)}
    {\(3\)}
    {\(4\)}
\end{problem}

\begin{answer}
D
\end{answer}

\begin{solution}
四个命题都不成立. 取 \(\alpha\) 为平面 \(z=0\).

对于第一个命题，取直线 \(m:(t,0,t+1)\) 和 \(n:(0,s,s+1)\). 它们在 \(\alpha\) 内的射影分别为方向向量为 \((1,0,0)\) 和 \((0,1,0)\) 的两条直线，二者垂直；但 \(m\)，\(n\) 的方向向量 \((1,0,1)\)，\((0,1,1)\) 的数量积为 \(1\)，所以 \(m\) 与 \(n\) 不垂直.

对于第二个命题，取 \(m:(t,0,t+1)\)，取方向向量为 \((1,1,-1)\) 且经过 \((0,0,1)\) 的直线 \(n:(s,s,1-s)\). 两条直线在 \((0,0,1)\) 相交，且方向向量数量积为 \(1-1=0\)，所以 \(m\perp n\)；但其射影方向向量 \((1,0,0)\) 与 \((1,1,0)\) 的数量积为 \(1\)，所以 \(m'\) 与 \(n'\) 不垂直.

对于第三个命题，取 \(m:(t,0,1)\)，\(n:(0,s,2)\). 射影 \(m'\) 和 \(n'\) 分别为两条相交的坐标轴，但 \(m\)，\(n\) 位于不同的水平面内，互不相交且不重合.

对于第四个命题，取 \(m:(t,0,t+1)\)，\(n:(s,1,2s+2)\). 射影 \(m'\)，\(n'\) 是两条平行直线，而 \(m\)，\(n\) 的方向向量分别为 \((1,0,1)\)，\((1,0,2)\)，不平行，且两直线不相交. 因此四个命题均不正确，不正确的命题个数为 \(4\)，故选 D.
\end{solution}

\begin{problem}
已知 \(p\) 和 \(q\) 是两个不相等的正整数，且 \(q \geqslant 2\)，则 \(\displaystyle\lim_{n\to\infty} \frac{\left(1 + \frac{1}{n}\right)^p - 1}{\left(1 + \frac{1}{n}\right)^q - 1}\) 等于（\(\quad\)）

\choices
    {\(0\)}
    {\(1\)}
    {\(\frac{p}{q}\)}
    {\(\frac{p - 1}{q - 1}\)}
\end{problem}

\begin{answer}
C
\end{answer}

\begin{solution}
对任意正整数 \(r\)，有
\[
\left(1+\frac{1}{n}\right)^r-1=\frac{1}{n}\sum_{j=0}^{r-1}\left(1+\frac{1}{n}\right)^j
\]
因此原极限等于
\[
\frac{\sum_{j=0}^{p-1}\left(1+\frac{1}{n}\right)^j}{\sum_{j=0}^{q-1}\left(1+\frac{1}{n}\right)^j}
\]
当 \(n\to\infty\) 时，分子趋于 \(p\)，分母趋于 \(q\)，故极限为 \(\dfrac{p}{q}\)，选 C.
\end{solution}

\begin{problem}
若数列 \(\{a_{n}\}\) 满足 \(\frac{a_{n+1}^2}{a_{n}^2} = p\)（\(p\) 为正常数，\(n \in \N^{*}\)），则称 \(\{a_n\}\) 为“等方比数列”. 甲：数列 \(\{a_n\}\) 是等方比数列；乙：数列 \(\{a_n\}\) 是等比数列，则（\(\quad\)）

\choices
    {甲是乙的充分条件但不是必要条件}
    {甲是乙的必要条件但不是充分条件}
    {甲是乙的充要条件}
    {甲既不是乙的充分条件也不是乙的必要条件}
\end{problem}

\begin{answer}
B
\end{answer}

\begin{solution}
若数列 \(\{a_n\}\) 是等比数列，设其公比为 \(r\). 由定义可知各项非零，且
\[
\frac{a_{n+1}^2}{a_n^2}=r^2
\]
所以它一定是等方比数列，即乙是甲的充分条件.

反过来，等方比数列的相邻两项绝对值之比相同，但相邻两项的正负号不一定按固定比例变化. 例如取 \(a_1=1\)，\(a_2=1\)，\(a_3=-1\)，并令其余各项为 \(1\)，则 \(a_{n+1}^2/a_n^2=1\)，但相邻项之比不是常数，故它不是等比数列. 所以甲是乙的必要条件但不是充分条件，选 B.
\end{solution}

\begin{problem}
双曲线 \(C_1: \frac{x^2}{a^2} - \frac{y^2}{b^2} = 1\)（\(a > 0\)，\(b > 0\)）的左准线为 \(l\)，左焦点和右焦点分别为 \(F_1\) 和 \(F_2\). 抛物线 \(C_2\) 的准线为 \(l\)，焦点为 \(F_2\)；\(C_1\) 与 \(C_2\) 的一个交点为 \(M\)，则 \(\frac{|F_1F_2|}{|MF_1|} - \frac{|MF_1|}{|MF_2|}\) 等于（\(\quad\)）

\choices
    {\(-1\)}
    {\(1\)}
    {\(-\frac{1}{2}\)}
    {\(\frac{1}{2}\)}
\end{problem}

\begin{answer}
A
\end{answer}

\begin{solution}
记双曲线的焦半距为 \(c=\sqrt{a^2+b^2}\)，则 \(F_1=(-c,0)\)，\(F_2=(c,0)\)，左准线为 \(x=-a^2/c\). 设交点 \(M=(x,y)\)，并记 \(d_1=|MF_1|\)，\(d_2=|MF_2|\). 由双曲线方程直接化简得
\(d_1^2=\left(\frac{c}{a}x+a\right)^2\)，\(d_2^2=\left(\frac{c}{a}x-a\right)^2\).
若 \(M\) 在左支，则 \(x\leqslant-a\). 因为 \(c>a\)，所以 \(x+a^2/c<0\)，且 \(\dfrac ca x-a<0\)，从而
\(d_2=-\left(\frac ca x-a\right)=a-\frac ca x\)，\(d_2=-\left(x+\frac{a^2}{c}\right)=-x-\frac{a^2}{c}\).
联立这两个等式得到
\[
x=\frac{a^2(c+a)}{c(c-a)}>0
\]
与 \(x\leqslant-a<0\) 矛盾，故 \(M\) 在右支. 此时
\(d_1=\frac{c}{a}x+a\)，\(d_2=\frac{c}{a}x-a\)，\(d_1-d_2=2a\).
又因 \(M\) 在以 \(F_2\) 为焦点、左准线为准线的抛物线上，\(d_2=x+a^2/c\). 联立
\[
\frac{c}{a}x-a=x+\frac{a^2}{c}
\]
得 \(d_2=\dfrac{2a^2}{c-a}\)，从而 \(d_1=\dfrac{2ac}{c-a}\). 于是
\[
\frac{|F_1F_2|}{|MF_1|}-\frac{|MF_1|}{|MF_2|}
=\frac{2c}{d_1}-\frac{d_1}{d_2}
=\frac{c-a}{a}-\frac{c}{a}=-1
\]
故选 A.
\end{solution}

\begin{problem}
已知两个等差数列 \(\{a_{n}\}\) 和 \(\{b_{n}\}\) 的前 \(n\) 项分别为 \(A_{n}\) 和 \(B_{n}\)，且 \(\frac{A_{n}}{B_{n}} = \frac{7n + 45}{n + 3}\)，则使得 \(\frac{a_{n}}{b_{n}}\) 为整数的正整数 \(n\) 的个数是（\(\quad\)）

\choices
    {\(2\)}
    {\(3\)}
    {\(4\)}
    {\(5\)}
\end{problem}

\begin{answer}
D
\end{answer}

\begin{solution}
由于等差数列前 \(n\) 项和是关于 \(n\) 的二次式且常数项为 \(0\)，设
\(A_n=\alpha n^2+\beta n\)，\(B_n=\gamma n^2+\delta n\).
由题设比例式，对所有正整数 \(n\) 有
\[
(\alpha n^2+\beta n)(n+3)=(\gamma n^2+\delta n)(7n+45)
\]
两边均为关于 \(n\) 的多项式，比较各次幂的系数，得
\(\alpha=7\gamma\)，\(3\alpha+\beta=45\gamma+7\delta\)，\(3\beta=45\delta\).
由此得到 \(\beta=15\delta\)，再代入中间的等式可得 \(\delta=3\gamma\)，从而 \(\alpha=7\gamma\)，\(\ \beta=45\gamma\). 因此
\(A_n=\gamma n(7n+45)\)，\(B_n=\gamma n(n+3)\)
其中 \(\gamma\neq0\).

令 \(A_0=B_0=0\)，相邻前项和之差给出
\(a_n=A_n-A_{n-1}=\gamma(14n+38)\)，\(b_n=B_n-B_{n-1}=\gamma(2n+2)\).
所以
\[
\frac{a_n}{b_n}=\frac{14n+38}{2n+2}=\frac{7n+19}{n+1}=7+\frac{12}{n+1}
\]
该数为整数当且仅当 \(n+1\mid12\). 正整数 \(n\) 对应 \(n+1=2\)，\(3\)，\(4\)，\(6\)，\(12\)，共有 \(5\) 个，故选 D.
\end{solution}

\begin{problem}
连掷两次骰子得到的点数分别为 \(m\) 和 \(n\)，记向量 \(\bs{a} = (m, n)\) 与向量 \(\bs{b} = (1, -1)\) 的夹角为 \(\theta\)，则 \(\theta \in \left(0, \frac{\pi}{2}\right]\) 的概率是（\(\quad\)）

\choices
    {\(\frac{5}{12}\)}
    {\(\frac{1}{2}\)}
    {\(\frac{7}{12}\)}
    {\(\frac{5}{6}\)}
\end{problem}

\begin{answer}
C
\end{answer}

\begin{solution}
向量 \(\bs a=(m,n)\) 与 \(\bs b=(1,-1)\) 的数量积为 \(m-n\). 夹角满足 \(0<\theta\leqslant\pi/2\) 时必须有 \(m-n\geqslant0\)，即 \(m\geqslant n\). 因为 \(m\)，\(n\) 均为 \(1\)，\(2\)，\(\ldots\)，\(6\)，且 \((m,n)\) 不可能是向量 \((1,-1)\) 的正倍数，所以 \(\theta>0\) 自动满足. 满足 \(m\geqslant n\) 的结果数为 \(6+5+4+3+2+1=21\)，总结果数为 \(36\)，故所求概率为 \(21/36=7/12\)，选 C.
\end{solution}

\begin{problem}
已知直线 \(\frac{x}{a} + \frac{y}{b} = 1\)（\(a\)，\(b\) 是非零常数）与圆 \(x^2 + y^2 = 100\) 有公共点，且公共点的横坐标均为整数，那么这样的直线共有（\(\quad\)）

\choices
    {\(60\) 条}
    {\(66\) 条}
    {\(72\) 条}
    {\(78\) 条}
\end{problem}

\begin{answer}
A
\end{answer}

\begin{solution}
按原试卷的字面条件，只要求公共点的横坐标为整数. 圆上横坐标为整数的点共有 \(40\) 个：当 \(x=\pm10\) 时各有一个点，当 \(x=-9\)，\(-8\)，\(\ldots\)，\(9\) 时各有两个点.

直线 \(\dfrac{x}{a}+\dfrac{y}{b}=1\) 的横、纵截距均非零，所以它不经过原点，也不平行于坐标轴. 由上述 \(40\) 个点任取两点确定的割线共有 \(\mathrm C_{40}^{2}\) 条. 其中竖直割线有 \(19\) 条（对应 \(x=-9\)，\(-8\)，\(\ldots\)，\(9\)），水平割线有 \(19\) 条（对应 \(y=\pm\sqrt{100-r^2}\)，\(r=1\)，\(2\)，\(\ldots\)，\(9\)，以及 \(y=0\)），经过原点的直径有 \(20\) 条. 竖直直径和水平直径分别与“经过原点”的情形重复一次，因此不合要求的割线共有
\[
19+19+20-2=56\text{ 条}
\]
所以符合条件的割线有
\[
\mathrm C_{40}^{2}-56=724\text{ 条}
\]

再看切线. 过上述 \(40\) 个点各有一条切线，其中 \(y=0\) 的两个点处切线为竖直线，\(x=0\) 的两个点处切线为水平线，均不能写成题设形式；其余 \(36\) 条切线符合题设. 因而按原试卷字面，直线总数为
\[
724+36=760\text{ 条}
\]
这不在选项中，故题面与选项及现有参考答案 A 不相容. 若将题干补为“公共点的横坐标和纵坐标均为整数”，才得到现有答案 \(60\) 条.
\end{solution}

\section{填空题}

\begin{problem}
已知函数 \( y = 2x - a \) 的反函数是 \( y = bx + 3 \)，则 \(a=\)\fillinblank{}；\(b=\)\fillinblank{}.
\end{problem}

\begin{answer}
\(6\)，\(\frac{1}{2}\)
\end{answer}

\begin{solution}
由 \(y=2x-a\) 解得 \(x=\dfrac{y+a}{2}\)，交换自变量与函数值后，反函数为 \(y=\dfrac{x+a}{2}\). 与已知反函数 \(y=bx+3\) 比较系数，得到 \(b=\dfrac12\)，且 \(a/2=3\)，所以 \(a=6\).
\end{solution}

\begin{problem}
复数 \(z = a + b\i\)，\(a\)，\(b \in \R\)，且 \(b \neq 0\). 若 \(z^2 - 4bz\) 是实数，则有序实数对 \((a, b)\) 可以是\fillinblank{}. （写出一个有序实数对即可）
\end{problem}

\begin{answer}
\((2,1)\)
\end{answer}

\begin{solution}
由 \(z=a+b\i\)，得
\[
z^2-4bz=(a^2-b^2-4ab)+(2ab-4b^2)\i
\]
它为实数，当且仅当虚部为零，即 \(2b(a-2b)=0\). 因为 \(b\neq0\)，所以 \(a=2b\). 取 \(b=1\)，则 \(a=2\)，故一个符合条件的有序实数对为 \((2,1)\).
\end{solution}

\begin{problem}
设变量 \(x\)，\(y\) 满足约束条件 \(\begin{cases}
x - y + 3 \geqslant 0,\\
x + y \geqslant 0,\\
-2 \leqslant x \leqslant 3,
\end{cases}\) 则目标函数 \(2x + y\) 的最小值为\fillinblank{}.
\end{problem}

\begin{answer}
\(-\frac{3}{2}\)
\end{answer}

\begin{solution}
由约束条件得 \(y\leqslant x+3\)，且 \(y\geqslant-x\). 要有可行的 \(y\)，必须满足 \(-x\leqslant x+3\)，即 \(x\geqslant-\dfrac32\). 目标函数为
\[
2x+y\geqslant2x-x=x
\]
又因 \(x\geqslant-\dfrac32\)，故 \(2x+y\geqslant-\dfrac32\). 当 \(x=-\dfrac32\)，\(y=\dfrac32\) 时，三个约束均满足，且目标函数等于 \(-\dfrac32\)，所以最小值为 \(-\dfrac32\).
\end{solution}

\begin{problem}
某男篮球运动员在三分线投球的命中率是 \(\frac{1}{2}\)，他投球 \(10\) 次，恰好投进 \(3\) 个球的概率为\fillinblank{}. （用数值作答）
\end{problem}

\begin{answer}
\(\frac{15}{128}\)
\end{answer}

\begin{solution}
设投进的次数为 \(X\). 各次投球相互独立，且投进、未投进的概率均为 \(1/2\)，所以 \(X\sim B\left(10,\dfrac12\right)\). 因而
\[
P(X=3)=\mathrm C_{10}^{3}\left(\frac12\right)^3\left(\frac12\right)^7=\frac{120}{2^{10}}=\frac{15}{128}
\]
\end{solution}

\begin{problem}
为了预防流感，某学校对教室用药熏消毒法进行消毒. 已知药物释放过程中，室内每立方米空气中的含药量 \(y\)（毫克）与时间 \(t\)（小时）成正比；药物释放完毕后，\(y\) 与 \(t\) 的函数关系式为 \(y = \left(\frac{1}{16}\right)^{t - a}\)（\(a\) 为常数）. 如图所示，据图中提供的信息，回答下列问题：

\begin{enumerate}
\item 从药物释放开始，每立方米空气中的含药量 \(y\)（毫克）与时间 \(t\)（小时）之间的函数关系式为\fillinblank{}；
\item 据测定，当空气中每立方米的含药量降低到 \(0.25\) 毫克以下时，学生方可进教室，那么药物释放开始，至少需要经过\fillinblank{} 小时后，学生才能回到教室.
\end{enumerate}

\bitmapfigure[max width=0.70\linewidth]{img/2007/hubei_science/q15_fig1.png}
\end{problem}

\begin{answer}
\(y=\begin{cases}
10t,&0\leqslant t\leqslant \frac{1}{10},\\
\left(\frac{1}{16}\right)^{t-\frac{1}{10}},&t>\frac{1}{10},
\end{cases}\)；\(0.6\)
\end{answer}

\begin{solution}
释放过程中 \(y\) 与 \(t\) 成正比，图象经过 \((0.1,1)\)，故此阶段的比例系数为 \(1/0.1=10\)，从而 \(y=10t\)，其中 \(0\leqslant t\leqslant0.1\). 释放结束后函数为 \(y=\left(\dfrac1{16}\right)^{t-a}\)，且图象在 \(t=0.1\) 处连续并取值 \(1\)，所以 \(\left(\dfrac1{16}\right)^{0.1-a}=1\)，得到 \(a=0.1\). 因而
\[
y=\begin{cases}
10t,&0\leqslant t\leqslant\dfrac1{10},\\
\left(\dfrac1{16}\right)^{t-\frac1{10}},&t>\dfrac1{10}.
\end{cases}
\]
当 \(t>0.1\) 时，要使含药量不超过 \(0.25=1/4=(1/16)^{1/2}\)，由于底数 \(1/16\) 小于 \(1\)，须有 \(t-0.1\geqslant1/2\)，即 \(t\geqslant0.6\). 所以至少经过 \(0.6\) 小时.
\end{solution}

\section{解答题}

\begin{problem}
已知 \(\triangle ABC\) 的面积为 \(3\)，且满足 \(0 \leqslant \overrightarrow{AB} \cdot \overrightarrow{AC} \leqslant 6\)，设 \(\overrightarrow{AB}\) 和 \(\overrightarrow{AC}\) 的夹角为 \(\theta\).

\begin{enumerate}
\item 求 \(\theta\) 的取值范围；
\item 求函数 \(f(\theta) = 2\sin^2\left(\frac{\pi}{4} + \theta\right) - \sqrt{3}\cos 2\theta\) 的最大值与最小值.
\end{enumerate}
\end{problem}

\begin{solution}
\begin{enumerate}
\item 记 \(u=|\overrightarrow{AB}|\,|\overrightarrow{AC}|\). 由三角形面积公式和数量积公式，得
\(u\sin\theta=6\)，\(0\leqslant u\cos\theta\leqslant6\).
因为 \(u>0\)，所以 \(0<\theta\leqslant\pi/2\). 再由 \(u=6/\sin\theta\)，得到 \(6\cot\theta\leqslant6\)，即 \(\cot\theta\leqslant1\). 在 \((0,\pi/2]\) 上余切函数递减，故 \(\theta\geqslant\pi/4\)，从而 \(\theta\in\left[\dfrac{\pi}{4},\dfrac{\pi}{2}\right]\).
\item 利用 \(2\sin^2\alpha=1-\cos2\alpha\)，有
\[
f(\theta)=1+\sin2\theta-\sqrt3\cos2\theta=1+2\sin\left(2\theta-\frac{\pi}{3}\right)
\]
当 \(\theta\in\left[\dfrac{\pi}{4},\dfrac{\pi}{2}\right]\) 时，\(2\theta-\dfrac\pi3\in\left[\dfrac\pi6,\dfrac{2\pi}3\right]\). 该区间内正弦函数的最大值为 \(1\)，在自变量为 \(\pi/2\) 时取得；最小值为 \(1/2\)，在左端点取得. 因此 \(f(\theta)\) 的最大值为 \(3\)，最小值为 \(2\).
\end{enumerate}
\end{solution}

\begin{problem}
在生产过程中，测得纤维产品的纤度（表示纤维粗细的一种量）共有 \(100\) 个数据，将数据分组如下表：

\begin{center}
\small
\begin{tabular}{|c|c|c|}
\hline
分组 & 频数 & 频率 \\
\hline
\([1.30,1.34)\) & \(4\) & \\
\hline
\([1.34,1.38)\) & \(25\) & \\
\hline
\([1.38,1.42)\) & \(30\) & \\
\hline
\([1.42,1.46)\) & \(29\) & \\
\hline
\([1.46,1.50)\) & \(10\) & \\
\hline
\([1.50,1.54)\) & \(2\) & \\
\hline
合计 & \(100\) & \\
\hline
\end{tabular}
\end{center}

\begin{enumerate}
\item 补全频率分布表，并画出频率分布直方图；
\item 估计纤度落在 \([1.38, 1.50)\) 中的概率及纤度小于 \(1.40\) 的概率分别是多少？
\item 统计方法中，同一组数据常用该组区间的中点值（例如区间 \([1.30, 1.34)\) 的中点值是 \(1.32\)）作为代表. 据此，估计纤度的期望.
\end{enumerate}

\bitmapfigure[max width=0.55\linewidth]{img/2007/hubei_science/q17_fig1.png}
\end{problem}

\begin{solution}
\begin{enumerate}
\item 各组频率等于频数除以总数 \(100\)，所以从第一组到第六组依次为 \(0.04\)，\(0.25\)，\(0.30\)，\(0.29\)，\(0.10\)，\(0.02\)，总和为 \(1.00\). 各组组距均为 \(0.04\)，故频率分布直方图中六个矩形的高，即频率/组距，依次为 \(1\)，\(6.25\)，\(7.5\)，\(7.25\)，\(2.5\)，\(0.5\)，底边分别为题目给出的六个区间.
\item 区间 \([1.38,1.50)\) 包含第三、四、五组，故其概率约为 \(0.30+0.29+0.10=0.69\). 对于小于 \(1.40\) 的概率，前两组贡献 \(0.04+0.25\)，第三组中 \([1.38,1.40)\) 占组距的一半，估计贡献为 \(0.30\times\dfrac{0.02}{0.04}=0.15\)，所以该概率约为 \(0.04+0.25+0.15=0.44\).
\item 各组中点依次为 \(1.32\)，\(1.36\)，\(1.40\)，\(1.44\)，\(1.48\)，\(1.52\)，用频率加权得到纤度期望的估计值
\[
1.32\times0.04+1.36\times0.25+1.40\times0.30+1.44\times0.29+1.48\times0.10+1.52\times0.02=1.4088
\]
\end{enumerate}
\end{solution}

\begin{problem}
如图，在三棱锥 \(V - ABC\) 中，\(VC \perp\) 底面 \(ABC\)，\(AC \perp BC\)，\(D\) 是 \(AB\) 的中点，且 \(AC = BC = a\)，\(\angle VDC = \theta\)（\(0 < \theta < \frac{\pi}{2}\)）.

\begin{enumerate}
\item 求证：平面 \(VAB \perp\) 平面 \(VCD\)；
\item 当角 \(\theta\) 变化时，求直线 \(BC\) 与平面 \(VAB\) 所成的角的取值范围.
\end{enumerate}

\bitmapfigure[max width=0.40\linewidth]{img/2007/hubei_science/q18_fig1.png}
\end{problem}

\begin{solution}
\begin{enumerate}
\item 因为 \(AC=BC\)，所以点 \(C\) 在 \(AB\) 的垂直平分面内；又因为 \(D\) 是 \(AB\) 的中点，所以 \(CD\perp AB\). 又因 \(VC\perp\) 底面 \(ABC\)，所以 \(VC\perp AB\). 直线 \(DC\)，\(VC\) 是平面 \(VCD\) 内相交的两条直线，且 \(AB\) 同时垂直于它们，故 \(AB\perp\) 平面 \(VCD\). 由于 \(AB\subset\) 平面 \(VAB\)，所以平面 \(VAB\perp\) 平面 \(VCD\).
\item 建立空间直角坐标系，取 \(C=(0,0,0)\)，\(A=(a,0,0)\)，\(B=(0,a,0)\)，\(V=(0,0,h)\)，其中 \(h>0\). 则 \(D=(a/2,a/2,0)\). 由 \(\angle VDC=\theta\)，在直角三角形 \(VCD\) 中有 \(CD=a/\sqrt2\)，故
\[
h=CD\tan\theta=\frac{a}{\sqrt2}\tan\theta
\]
平面 \(VAB\) 的一个法向量可取 \(\bs n=(h,h,a)\)，直线 \(BC\) 的方向向量可取 \(\bs d=(0,1,0)\). 设直线与平面所成的角为 \(\varphi\)，则
\[
\sin\varphi=\frac{|\bs d\cdot\bs n|}{|\bs d|\,|\bs n|}=\frac{h}{\sqrt{2h^2+a^2}}=\frac{\sin\theta}{\sqrt2}
\]
因为 \(0<\theta<\pi/2\)，所以 \(0<\sin\varphi<1/\sqrt2\)，从而 \(0<\varphi<\pi/4\). 当 \(\theta\) 在给定范围内变化时，\(\sin\theta\) 可取所有 \((0,1)\) 内的值，故取值范围为 \(\left(0,\dfrac\pi4\right)\).
\end{enumerate}
\end{solution}

\begin{problem}
在平面直角坐标系 \(xOy\) 中，过定点 \(C(0,p)\) 作直线与抛物线 \(x^{2} = 2py\)（\(p > 0\)）相交于 \(A\)，\(B\) 两点.

\begin{enumerate}
\item 若点 \(N\) 是点 \(C\) 关于坐标原点 \(O\) 的对称点，求 \(\triangle ANB\) 面积的最小值；
\item 是否存在垂直于 \(y\) 轴的直线 \(l\)，使得 \(l\) 被以 \(AC\) 为直径的圆截得的弦长恒为定值？若存在，求出 \(l\) 的方程；若不存在，说明理由.
\end{enumerate}

\FigureLayoutDeclare{img/2007/hubei_science/q19_fig1.png}{10.0cm}{22bp 36bp 1bp 19bp}
\bitmapfigure[max width=0.60\linewidth]{img/2007/hubei_science/q19_fig1.png}
\end{problem}

\begin{solution}
\begin{enumerate}
\item 设过 \(C(0,p)\) 的直线为 \(y=kx+p\)，与抛物线联立得
\[
x^2-2pkx-2p^2=0
\]
设两根为 \(x_1\)，\(x_2\)，则 \(x_1+x_2=2pk\)，\(x_1x_2=-2p^2\). 交点为 \(A=(x_1,kx_1+p)\)，\(B=(x_2,kx_2+p)\)，而 \(N=(0,-p)\). 因此
\[
2S_{\triangle ANB}=\left|\begin{vmatrix}x_1&kx_1+2p\\x_2&kx_2+2p\end{vmatrix}\right|=2p|x_1-x_2|
\]
由根的判别式，\(|x_1-x_2|=2p\sqrt{k^2+2}\)，所以
\[
S_{\triangle ANB}=2p^2\sqrt{k^2+2}\geqslant2\sqrt2p^2
\]
当 \(k=0\) 时取得等号，故面积的最小值为 \(2\sqrt2p^2\).
\item 设抛物线上的点 \(A\) 为 \(\left(u,\dfrac{u^2}{2p}\right)\)，则以 \(AC\) 为直径的圆上点 \((X,Y)\) 满足
\[
(X-u)X+\left(Y-\frac{u^2}{2p}\right)(Y-p)=0
\]
若水平直线为 \(Y=t\)，代入圆方程得关于 \(X\) 的二次方程，其两根之差的平方，也就是弦长平方，为
\[
L^2=u^2-4\left(t-\frac{u^2}{2p}\right)(t-p)=u^2\left(\frac{2t}{p}-1\right)+4t(p-t)
\]
要使弦长与任意的 \(A\)，即与变化的 \(u\)，无关，必须有 \(2t/p-1=0\)，所以 \(t=p/2\). 此时 \(L^2=p^2\)，弦长恒为 \(p\)，故存在所求直线，方程为 \(l:y=\dfrac p2\).
\end{enumerate}
\end{solution}

\begin{problem}
已知定义在正实数集上的函数 \(f(x) = \frac{1}{2}x^2 + 2ax\)，\(g(x) = 3a^2\ln x + b\)，其中 \(a > 0\). 设两曲线 \(y = f(x)\)，\(y = g(x)\) 有公共点，且在该点处的切线相同.

\begin{enumerate}
\item 用 \(a\) 表示 \(b\)，并求 \(b\) 的最大值；
\item 求证：\(f(x) \geqslant g(x)\)（\(x > 0\)）.
\end{enumerate}
\end{problem}

\begin{solution}
\begin{enumerate}
\item 设公共点的横坐标为 \(x_0>0\). 由两曲线在该点相切，得
\[
x_0+2a=f'(x_0)=g'(x_0)=\frac{3a^2}{x_0}
\]
于是 \(x_0^2+2ax_0-3a^2=(x_0-a)(x_0+3a)=0\). 因为 \(x_0>0\)，\(a>0\)，所以 \(x_0=a\). 再由公共点的函数值相等，得
\(\frac52a^2=3a^2\ln a+b\)，\(b=\frac52a^2-3a^2\ln a\).
令 \(B(a)=\frac52a^2-3a^2\ln a\)，则
\[
B'(a)=2a-6a\ln a=2a(1-3\ln a)
\]
故 \(B\) 在 \((0,\e^{1/3})\) 上递增，在 \((\e^{1/3},+\infty)\) 上递减；又 \(B(a)\to0\)（\(a\to0^+\)），且 \(B(a)\to-\infty\)（\(a\to+\infty\)），所以最大值在 \(a=\e^{1/3}\) 处取得，为 \(\dfrac32\e^{2/3}\).
\item 令 \(h(x)=f(x)-g(x)\)，代入上式的 \(b\) 得 \(h(a)=0\). 且
\[
h'(x)=x+2a-\frac{3a^2}{x}=\frac{(x-a)(x+3a)}{x}
\]
当 \(0<x<a\) 时，\(h'(x)<0\)；当 \(x>a\) 时，\(h'(x)>0\). 所以 \(h\) 在 \(x=a\) 处取得全局最小值 \(0\)，即对任意 \(x>0\)，都有 \(f(x)-g(x)\geqslant0\)，从而 \(f(x)\geqslant g(x)\).
\end{enumerate}
\end{solution}

\begin{problem}
已知 \(m\)，\(n\) 为正整数.

\begin{enumerate}
\item 用数学归纳法证明：当 \(x > -1\) 时，\((1 + x)^{m} \geqslant 1 + mx\)；
\item 对于 \(n \geqslant 6\)，已知 \(\left(1 - \frac{1}{n + 3}\right)^n < \frac{1}{2}\). 求证 \(\left(1 - \frac{m}{n + 3}\right)^m < \frac{1}{2}\)，\(m = 1\)，\(2\)，\(\cdots\)，\(n\)；
\item 求满足等式 \(3^{n} + 4^{n} + \cdots + (n + 2)^{n} = (n + 3)^{n}\) 的所有正整数 \(n\).
\end{enumerate}
\end{problem}

\begin{solution}
\begin{enumerate}
\item 当 \(m=1\) 时，\((1+x)^1=1+x\)，不等式成立. 假设对某个正整数 \(m\) 成立，即 \((1+x)^m\geqslant1+mx\). 因为 \(x>-1\)，所以 \(1+x>0\)，从而
\[
(1+x)^{m+1}\geqslant(1+mx)(1+x)=1+(m+1)x+mx^2\geqslant1+(m+1)x
\]
由数学归纳法，结论对一切正整数 \(m\) 成立.
\item 这一小问按原试卷的字面式子不成立. 取 \(n=6\)，则前提为 \(\left(1-\dfrac19\right)^6=(8/9)^6<1/2\)，因为 \(2\cdot8^6=524288<531441=9^6\). 但取允许的 \(m=1\) 时，结论给出 \(\left(1-\dfrac19\right)^1=8/9<1/2\)，这是错误的. 因此原题第（2）小问存在印刷或题面错误，不能对该命题给出合法证明；现稿原先的“结论成立”不保留.
\item 直接讨论原等式. 对 \(n=1\)，\(2\)，\(3\)，\(4\)，\(5\)，左端分别为 \(3\)，\(25\)，\(216\)，\(2258\)，\(28975\)，右端分别为 \(4\)，\(25\)，\(216\)，\(2401\)，\(32768\)，所以 \(n=2\)，\(3\) 成立，而 \(n=1\)，\(4\)，\(5\) 不成立.

当 \(n\geqslant6\) 时，令 \(r_n=\left(1-\dfrac1{n+3}\right)^n\). 对实数 \(x\geqslant6\)，令 \(H(x)=x\ln\left(1-\dfrac1{x+3}\right)\)，则
\[
H'(x)=\ln\left(1-\frac1{x+3}\right)+\frac{x}{(x+2)(x+3)}< -\frac1{x+3}+\frac{x}{(x+2)(x+3)}<0
\]
所以 \(r_n\) 递减，且 \(r_n\leqslant r_6=(8/9)^6<1/2\). 令 \(k=n+3\)，则
\[
\frac{3^n+4^n+\cdots+(n+2)^n}{(n+3)^n}=\sum_{j=1}^{n}\left(1-\frac{j}{k}\right)^n
\]
由第（1）问对 \(x=-1/k\) 的结论，\(\left(1-1/k\right)^j\geqslant1-j/k>0\)，故
\[
\left(1-\frac{j}{k}\right)^n\leqslant\left(1-\frac1k\right)^{jn}=r_n^j
\]
于是
\[
\frac{3^n+4^n+\cdots+(n+2)^n}{(n+3)^n}<\sum_{j=1}^{\infty}r_n^j=\frac{r_n}{1-r_n}<1
\]
因此 \(n\geqslant6\) 时等式不可能成立，综上所有正整数解为 \(n=2\)，\(3\).
\end{enumerate}
\end{solution}
